\chapter{Black Eye}

\section*{Definition and Overview}
A black eye, also called a periorbital haematoma, is bruising of the skin and soft tissues around the eye, usually caused by a blow to the face. Blood collects under the thin skin of the eyelid and eye socket, producing the typical purple, blue, or black discolouration, which later fades to green and yellow as it resolves.

Most black eyes are minor injuries that settle on their own within one to two weeks. However, any injury that affects the eyeball itself, the vision, or the bones surrounding the eye requires urgent medical assessment, because damage to these structures can be sight-threatening.

\section*{Epidemiology}
Black eyes are very common and can occur at any age.

\begin{itemize}
    \item \textbf{Common causes:} accidental falls, sporting injuries, workplace accidents, and fights or assaults
    \item \textbf{Emergency attendance:} a frequent reason for visits to emergency departments, especially among children, young adults, and older people taking blood-thinning medicines
    \item \textbf{Associated injuries:} most are simple soft-tissue bruises, but a minority are associated with injury to the eyeball or fractures of the facial bones
    \item \textbf{Risk factors:} contact sports, hazardous work, and anticoagulant therapy increase both the risk and severity of the injury
\end{itemize}

\section*{Aetiology and Pathophysiology}
A black eye occurs when blunt force damages the small blood vessels beneath the thin skin around the eye.

\begin{itemize}
    \item \textbf{Direct trauma:} a fist, ball, branch, or fall strikes the eye area
    \item \textbf{Bleeding under the skin:} damaged vessels leak blood into the soft tissues around the orbit, causing swelling and discolouration
    \item \textbf{Accumulation and settling:} gravity draws the bruising downward, so discolouration may spread onto the cheek or temple
    \item \textbf{Blood-thinning medicines:} aspirin, warfarin, and other anticoagulants increase bleeding and can make bruising far more extensive
    \item \textbf{Bleeding disorders:} conditions that impair clotting can cause prominent bruising after only minor injury
    \item \textbf{Associated injuries:} a sufficiently forceful blow can also damage the eyeball, the fluid within it, or the thin bones of the eye socket
\end{itemize}

\section*{Clinical Features}
\begin{itemize}
    \item Swelling and bruising around the eye, which may be black, blue, or purple at first
    \item Pain and tenderness over the affected area
    \item Discolouration spreading to the cheek or temple over one to two days
    \item Bruising changing colour to green and yellow as it heals
    \item \textbf{Warning features:} double vision, blurred or reduced vision, severe pain, an inability to move the eye fully, or numbness of the cheek
    \item A step or unevenness felt in the bone around the eye, which suggests a fracture
    \item Persistent nosebleed or blood in the white of the eye, which may accompany a facial fracture
\end{itemize}

\section*{Differential Diagnosis}
\begin{itemize}
    \item Simple bruising of the eyelid and eye socket from minor trauma
    \item A black eye from a blow, which must be distinguished from bleeding caused by another condition
    \item Fracture of the floor or rim of the eye socket (blowout fracture)
    \item Injury to the eyeball itself, including hyphaema (blood inside the eye) or damage to the retina
    \item Raised pressure inside the eye after trauma
    \item Raccoon eyes, bruising around both eyes after a head injury, and other signs of a skull base fracture
    \item Allergic or contact dermatitis, or a severe eyelid infection such as periorbital cellulitis, which can mimic bruising and swelling
    \item Bleeding disorders or blood-thinning medicines as a cause of disproportionate bruising
\end{itemize}

\section*{When to Seek Medical Care}
Seek urgent medical care for a black eye with reduced or double vision, severe pain, an eye that cannot move normally, bleeding inside the eye, or signs of a facial fracture. Anyone taking blood-thinning medicines, and any child who has sustained a head injury, should be assessed early.

\textbf{Call 000 (or local emergency services) for:}
\begin{itemize}
    \item Loss of vision or a major change in vision after an eye injury
    \item A black eye together with confusion, vomiting, or loss of consciousness, which may indicate a head injury
    \item Uncontrolled bleeding from the eye or nose
    \item Intense eye pain with the eye feeling hard and the vision dim, which may indicate raised pressure inside the eye
\end{itemize}

\section*{Diagnosis and Investigations}
\begin{itemize}
    \item \textbf{History:} how the injury happened, its timing, whether vision is affected, and any use of blood-thinning medicines
    \item \textbf{Examination of the eye:} checking vision, pupil reactions, eye movements, and the front of the eye with a light
    \item \textbf{Eye pressure measurement:} to detect raised pressure after injury
    \item \textbf{Slit-lamp examination:} a microscope examination of the eye to look for hyphaema and other injuries
    \item \textbf{Imaging:} X-ray or CT scan of the face when a fracture is suspected
    \item \textbf{Assessment of the cheek:} checking for numbness, which may indicate nerve injury accompanying a fracture
    \item \textbf{Referral:} to an eye specialist (ophthalmologist) or facial surgeon if significant injury is found
\end{itemize}

\section*{Management}
\begin{itemize}
    \item \textbf{Cold packs:} apply a wrapped cold pack or ice pack for 15 to 20 minutes at a time during the first 24 to 48 hours to reduce swelling
    \item \textbf{Simple pain relief:} paracetamol if needed; avoid aspirin and anti-inflammatory medicines for the first day or two, because they can increase bleeding
    \item \textbf{Rest:} avoid further contact to the area and heavy exertion
    \item \textbf{Elevation:} keeping the head slightly raised when sleeping can reduce swelling
    \item \textbf{Warm compresses:} after the first two days, gentle warmth may help the bruise clear
    \item \textbf{Surgery:} needed only if there is a displaced facial fracture or a trapped eye muscle
    \item \textbf{Review:} a follow-up eye check if vision or eye-movement symptoms are present
\end{itemize}

\section*{Complications}
\begin{itemize}
    \item \textbf{Hyphaema:} blood pooling in the front chamber of the eye, which requires urgent eye care
    \item \textbf{Raised pressure inside the eye:} which can damage the optic nerve
    \item \textbf{Blowout fracture of the eye socket:} with trapped eye muscles causing persistent double vision
    \item \textbf{Damage to the retina or lens:} caused by the force of the blow
    \item \textbf{Infection:} of the bruised tissue or of blood within the eye, in rare cases
    \item \textbf{Facial asymmetry or scarring:} if fractures are not treated
    \item \textbf{Cosmetic and psychological impact:} particularly when the injury resulted from an assault
\end{itemize}

\section*{Prognosis and Outcomes}
Most black eyes are minor soft-tissue bruises that resolve fully within one to two weeks as the trapped blood is absorbed. Swelling usually settles within a few days, and the bruise fades from black and blue through green and yellow, leaving no permanent mark.

Outcomes are excellent when the eyeball and bones are unharmed. Injuries involving the eyeball, vision, or an orbital fracture require specialist treatment and may take considerably longer to heal, and some carry a risk of lasting visual impairment.

\section*{Prevention and Self-Care}
\begin{itemize}
    \item Wear appropriate eye and face protection for sports such as squash, cricket, and paintball
    \item Use recommended protective equipment, including helmets with face guards, for high-risk work and sport
    \item Secure or remove hazards around the home that can cause falls
    \item Supervise children during active play
    \item Avoid physical confrontation that could involve a blow to the face
    \item If you take blood-thinning medicines, report any head or face injury promptly, even if it seems minor
    \item Apply a cold pack early after any blow to the face to limit swelling and bruising
\end{itemize}

\section*{Sources and Further Reading}
The following organisations provide reliable information on black eyes and eye injuries for patients, families, and health professionals.

\begin{itemize}
    \item NHS. \textit{Information on black eyes and when to seek medical help.} https://www.nhs.uk/
    \item Mayo Clinic. \textit{Overview of eye injuries, including first aid and treatment.} https://www.mayoclinic.org/
    \item MedlinePlus. \textit{Health information on black eye and eye injuries.} https://medlineplus.gov/
    \item MSD Manual. \textit{Professional information on eye trauma and periorbital injury.} https://www.msdmanuals.com/
    \item Centers for Disease Control and Prevention. \textit{Resources on injury prevention and head and eye safety.} https://www.cdc.gov/
\end{itemize}
